\documentclass[script,con]{debate}
\motion{语言的边界是 / 不是人类的边界}
\stance{反方}{不是}
\meta{赛制}{招新赛 3v3}
\meta{口径来源}{备赛文档 第十节 反方口径表}
\begin{document}
\debatetitle
\begin{thesisbox}[全场一句话立论]语言记下的是人走过的路。人总是先想到、先做到，词才跟上来；所以人能走到的地方，比语言能说到的地方远。\end{thesisbox}

\begin{kit}{1. 反一开篇立论}{套件 · 组装后 780–785 字 / 180 秒}
\assembly{开头（任选 1，接住刚结束的质询） → 立论主体}
\begin{pick}{开头}{1}
\begin{module}{如果质询拿到了对方承认失语的人还能思考}
谢谢主席。刚才的质询里，对方承认了一件事：失语的人还能思考。这句话，就是今天的答案。
\end{module}
\begin{module}{如果质询里对方回避了失语的人能不能思考}
谢谢主席。失语的人还能不能思考，刚才对方没有正面回答。没关系，我方先回答：能。
\end{module}
\begin{module}{如果质询拿到了对方承认比划也教得会}
谢谢主席。刚才的质询里，对方承认了：不出声，只比划，也教得会。不用语言，人也能把本事传下去。
\end{module}
\begin{fallback}{质询没有拿到明确成果}
谢谢主席。维特根斯坦那句名言，说的是“我的”语言、“我的”世界。今天的辩题问的是人类。
\end{fallback}
\end{pick}
\begin{fixed}{立论主体}
先说清三个词。语言，是有词、有语法、用来表达和交流意思的符号系统，口语、文字、手语都是；音乐、绘画叫语言，只是比喻。人类的边界，是人能想到、认识到、做到，并且能传给别人的范围的极限。“是”，指两条边界基本重合。

所以判断本题的标准是：语言到不了的地方，人类还能不能稳定地到达。反方要证明，在语言之外，人能稳定地思考、做到，而且留得住；正方要证明，人超出本能的那一截以语言为条件，语言之外的，要么到不了，要么留不住。只证明语言很有用，不够。

我方第一点：想在说前面。

思考不是用语言做的。大脑里有一套专门处理语言的网络。你听懂一句话，它就忙起来；可你算数、解逻辑题、听音乐的时候，它基本不参与。

所以，语言坏了，思想还在。二〇〇五年美国科学院院刊报道了三个人，左脑大面积损伤毁掉了他们的语法，说话连不成完整的句子，可带括号的算式，他们照样解。研究全面失语的学者总结：这些人几乎完全不能理解和说出语言，却还能加减，能解逻辑题，能推断别人的想法，能欣赏音乐。苏联作曲家舍巴林，中风后严重失语，还在作曲，作品被评为和病前一样好。

二〇二四年《自然》上一篇综述的结论是：语言不像是复杂思考的前提。它是说给别人听的工具，不是用来想的材料。

我方第二点：知道的比说的多。

你能在车站一千张脸里，一眼认出你妈。可你说得出她的鼻子多高、眼睛多宽吗？说不出。哲学家波兰尼把这叫默会知识：我们知道的，比我们能说出来的多。

人最高处的本事，都是这样。棋手的棋感，医生看一眼就觉得不对劲的直觉，老师傅手上的分寸。说不清，却一再做得到，还教得会。二〇一五年有个实验，教人打石器，只许比划、不许出声的那一组，每一下敲出好石片的概率，翻了一倍。我方承认，能说话的那组更快，翻了四倍。可快，是工具好用。汽车比走路快，公路修到的地方，不是人类的边界。

所以我方全场会请对方回答一个问题：\stress{一个中风失语的人，还能不能思考？}如果能，他的语言缩到几乎为零，他的世界没有跟着缩。那语言的边界，就不是人的边界。

说到底，语言是人类最好的工具，它记下人走过的路。可人总是先想到、先做到，词才跟上来。谢谢。
\end{fixed}
\end{kit}
\begin{contingency}
\ifthen{对方说失语者都是先有语言的}{孩子身上两套系统早就分开；没有语言的聋孩子照样学会数学、做因果推理，只是有些慢。慢不是边界。}
\ifthen{对方把语言定义成一切符号}{那题目就成了“人能表达的边界是人能表达的边界”，请回到有词有语法的中立定义。}
\ifthen{对方用轮扁后半句}{“不能以喻”是用话教不会；轮扁七十岁还在做，人到得了。}
\ifthen{对方追“说话翻四倍”}{大方承认更快，再说一句：快是工具好用，不是边界。}
\end{contingency}

\begin{stage}{2. 反一被质询防守预案}{正二质询，90 秒双边计时}
\textbf{原则}：先答是或否，再用一句话守住前提。承认语言有用、说话更快；守住的只有一句：工具不是边界。对方打断时停下来，每答不超过十秒。
\begin{qa}
\qarow{您方举的失语者，发病前会不会说话？}{会。可孩子身上，语言和思考早就是两套系统；没有语言的聋孩子，照样学会数学，只是有些慢。}
\qarow{他们的算术，是不是在会说话的时候学会的？}{学的时候有语言；算的时候不靠语言。题目问的是边界，不是学历。}
\qarow{一辈子没语言的人，有能精确数过三的例子吗？}{皮拉罕人没有数词，东西摆在眼前，很多个也一个不差地配对。数词是帮人记住的工具。}
\qarow{说不出的那一点，是不是每一代都要重新摸？}{是。每一代都有人摸到，这正说明人到得了。}
\qarow{师傅走了，徒弟没摸到，人类还有没有这一点？}{在下一个摸到的人手里。手艺几千年没断，靠的就是一代代人去摸。}
\qarow{比划翻一倍，说话翻四倍，差距是不是很大？}{是，说话更快。快是工具好用，不是边界。}
\qarow{维果茨基说思想是在词里完成的？}{同一章他也说，思想和词不是一个模子刻出来的。}
\qarow{轮扁说古之人与其不可传也死矣？}{他说的“不能以喻”，是用话教不会。他七十岁还在斫轮，人到得了。}
\end{qa}
\end{stage}

\begin{stage}{3. 反二质询正一问题链}{90 秒，双边计时}
\textbf{总目标}：全场第一次质询，拿到两个承认，交给反一立论开头用：失语的人还能思考；不用语言也教得会。问题要短，对方绕开时先复述一次，第二次再礼貌打断。
\begin{chain}{链一：失语}{失语的人还能思考}
\q{中风失语的人，还能不能算数？}{答“能”：下一问；答“不能”：请问出处，研究说能}
\q{他还能不能认出自己的孩子？}{答“能”：下一问}
\q{那他的世界，有没有跟着他的语言一起变小？}{答“没有”：收束；答“有”：请问小在哪里}
\cut{好，我换个问法：他还能不能想事情？}
\close{谢谢。对方承认，失语的人还能思考。}
\end{chain}
\begin{chain}{链二：比划}{不用语言也教得会}
\q{石器实验里，只比划的那组有没有提高？}{答“有”：下一问}
\q{比划算不算您方定义里的语言？}{答“不算”：下一问；答“算”：请回到双方的中立定义}
\q{那不用语言，是不是也教得会？}{答“是”：收束；答“只是慢”：慢不是到不了}
\close{谢谢。对方承认，不用语言也教得会。}
\end{chain}
\end{stage}

\begin{stage}{4. 反三对辩要点}{双方各 90 秒}
\textbf{目标}：开第一个战场，思考。让评委看到：语言坏了，思想还在；语言没到，人已经先到。每次发言先一句回应，再抛一个问题。
\begin{points}{进攻点（4–5 个，每个 15–20 秒）}
\item 中风失语的人还能不能思考？能的话，他的语言缩到几乎为零，他的世界缩了吗？
\item 尼加拉瓜那门手语是谁造的？造出来之前，他们是不是在语言外面？
\item 五个月的婴儿能算一加一，那时候他有语言吗？
\item 爱因斯坦说词语要到第二阶段才去找，那第一阶段，他在用什么想？
\item 皮拉罕人东西摆在眼前，一个不差地配对。他们缺的是数，还是记数的工具？
\end{points}
\begin{points}{备用点（6–8 个）}
\item 对方说失语者先有语言：孩子身上两套系统早就分开，语言剥夺的聋孩子照样学会数学。
\item 对方说二十五年替代不了：替代不了的是那套词的教学，而那套词，是一群没有语言的孩子造出来的。
\item 对方说婴儿的一加一动物也会：那就说明思考在语言之前就有，边界不在语言那里。
\item 对方说一个人摸到不算人类走到：人类就是一个个人，一个人到了，人类的边界就在那里。
\item 对方说维果茨基：同一章他说思想和词不是一个模子刻出来的。
\item 对方说说话翻四倍：快四倍是工具好用；汽车快，公路不是边界。
\item 对方说言语干扰会拖慢俄国人辨色：慢一百多毫秒，照样分得清。
\end{points}
\begin{points}{对方最可能问的三个问题 → 回应}
\item 失语者是不是先有语言的？→ 是，可算的时候不靠语言；孩子身上两套系统早就分开。
\item 说不出的东西，怎么让下一个人接着往前走？→ 做给他看，手把手，比划也翻一倍；每一代都有人到得了。
\item 一辈子没有语言的人，能不能理解别人想错了？→ 会慢。慢，不是到不了。
\end{points}
\end{stage}

\begin{kit}{5. 反二申论}{套件 · 组装后 380–398 字 / 90 秒}
\assembly{开头 → 拆正方立论（任选 1） → 推进论点一 → 收尾}
\begin{fixed}{开头}
谢谢主席。我先说这一篇做什么：先拆开正方的立论，再往前推我方的第一点。
\end{fixed}
\begin{pick}{拆正方立论}{1}
\begin{module}{如果对方立论建立在尼加拉瓜手语者答不对错误信念题上}
正方立论的主干，是尼加拉瓜第一批手语者答不对那道题。我方先认：学到那些词，确实帮他们学会了。可这证明的是词能帮人学会，不是语言到不了的地方，人就到不了。你想想，这门手语是谁造的？是一群原本没有语言的聋孩子。他们站在语言外面，自己走出了一门语言。人先到了，词才出来。何况第一批只有八个人；《自然》那篇综述也写，没有语言的人有些推理会慢，可照样学得会数学。慢，不是边界。
\end{module}
\begin{module}{如果对方立论建立在石器实验说话翻四倍上}
正方立论的主干，是石器实验里说话翻了四倍。我方承认，说话快。可同一个实验里，只许比划、不许出声的那组，也翻了一倍：不用语言，照样教得会。再看那七十多万年：没有像样的语言，石器照样一代一代传了七十多万年，一代也没断。快四倍，是工具好用。汽车比走路快得多，可公路修到的地方，从来不是人类的边界；公路外面，照样有人走路，有人翻山，有人第一个走到。这个实验，恰恰证明了人在语言之外也走得动。
\end{module}
\begin{module}{如果对方立论建立在维特根斯坦那句名言上}
正方立论的起点，是维特根斯坦那句名言。可原书里，那句话说的是“我的”语言、“我的”世界；几段之后他就写，哲学上的“我”，不是人，不是人的身体。更要紧的是，同一本书快结尾的地方他写：确实有不可说的东西，它显示自身。他给出版人写信说，这本书分两部分，写出来的，和没写出来的，没写出来的那部分才重要。所以正方拿来开场的那句话，出处本身就站在我方这一边。名言替不了论证。
\end{module}
\begin{fallback}{对方立论的主干不在以上几条时}
我方请评委回到判准。正方要证明的，不是语言很有用，而是语言到不了的地方，人就稳定地到不了。一个例子说明语言帮了忙，证明不了两条边界重合。反过来，只要找得到一块人稳定到达、语言却说不出的地方，“是”字就不成立。这样的地方，我方立论已经给了两块：一块是失语的人还在思考，一块是说不出的本事还在一代代传。请评委带着这两块地方，听后面的比赛。
\end{fallback}
\end{pick}
\begin{fixed}{推进论点一}
我再往前推一步。人不光在语言坏了以后还能想，还能在语言到之前就先想到。数学家阿达马写书之前，问了很多科学家：你们是怎么想问题的？爱因斯坦回了一封信，收在他一九四五年出的书里。他说，写出来或说出来的词语，在他的思考里似乎不起任何作用；他想的时候，用的是可以自由组合的图像和符号。词语，要到第二阶段，才费力地去找。说白了，他先到了那里，词是后来才去接他的。
\end{fixed}
\begin{fixed}{收尾}
所以请评委记住我方全场的问题，也请正方下一轮正面回答：一个中风失语的人，还能不能思考？能，还是不能？
\end{fixed}
\end{kit}

\begin{stage}{6. 反二被质询防守预案}{正三质询，90 秒双边计时}
\textbf{原则}：爱因斯坦的信要报得出出处：阿达马《数学领域的发明心理学》，一九四五年，附录二。对方一定会追“第二阶段是为了告诉别人”，先认，再守“先到的是人”。
\begin{qa}
\qarow{爱因斯坦说，词要到第二阶段才去找，对吗？}{对。}
\qarow{找词，是为了能告诉别人，对吗？}{对，为了告诉别人，不是为了想到。想到在前面。}
\qarow{要是他一直没找到，人类知不知道相对论？}{会晚一点知道。可那一步是他先走出去的；词是去接他的，不是先去划他的。}
\qarow{奥杜威石器七十多万年有没有明显变好？}{没怎么变。可它一代代传了七十多万年，没有像样的语言，人照样守得住。}
\qarow{人类的边界，是守住的地方，还是走出去的地方？}{都是。走出去的第一步，常常在词前面，爱因斯坦就是。}
\qarow{尼加拉瓜第一批只有八个人，你方凭什么说样本小？}{这是研究自己报告的人数；我方不否认结果，只说它证明的是“学会”，不是“边界”。}
\end{qa}
\end{stage}

\begin{stage}{7. 反三质询正二问题链}{90 秒，双边计时}
\textbf{总目标}：拿到两个承认：尼加拉瓜手语是没有语言的人造出来的；家庭手势者缺的是记大数的工具，不是数量感。
\begin{chain}{链一：造语言}{人先到，词后到}
\q{尼加拉瓜手语，是不是那群聋孩子自己造出来的？}{答“是”：下一问}
\q{造之前，他们有没有这门语言？}{答“没有”：下一问}
\q{那造出语言的人，是不是先走到了语言外面？}{答“是”：收束；答“不是”：请问语言从哪里来}
\close{谢谢。对方承认：造出语言的，是一群还没有语言的人。}
\end{chain}
\begin{chain}{链二：手势者}{缺的是工具}
\q{家庭手势者会不会用手势表达数量？}{答“会”：下一问}
\q{一到三个东西，他们是不是全对？}{答“是”：下一问}
\q{那他们缺的是数量感，还是记住大数的工具？}{答“工具”：收束；答“数量感”：他们数得接近，研究原话}
\cut{好，我换个问法：他们数得接近，对不对？}
\close{谢谢。缺的是工具，不是人到不了。}
\end{chain}
\end{stage}

\begin{stage}{8. 反一对辩要点}{双方各 90 秒}
\textbf{目标}：收拢前两段，把全场分歧钉在“工具还是边界”。正方一定会主打“一个人摸到不等于人类走到”，一辩要把这句话拆开。
\begin{points}{进攻点（4–5 个，每个 15–20 秒）}
\item 一个人先到了，这一步算不算人类的？不算的话，人类的边界是谁的边界？
\item 爱因斯坦想到相对论那天，人类有没有多一点东西？
\item 按正方的标准，没写进书里的手艺，算不算人类的？几千年的手艺人算不算人类？
\item 失语的人还能不能思考？他的世界，跟着他的语言缩了吗？
\end{points}
\begin{points}{备用点（6–8 个）}
\item 对方说说不出的传不下：比划组翻一倍，手艺几千年没断。
\item 对方说七十万年停滞：守得住，就是人到得了；停在那里的不是人，是速度。
\item 对方说性骚扰这个词：一九七五年以前她们就说得清，缺的是名字，名字是她们自己造的。
\item 对方说轮扁：七十岁还在斫轮，人到得了。
\item 对方说皮拉罕记不住：记不住是工具没带，不是人到不了。
\item 对方说脚手架：孩子身上两套系统早就分开。
\end{points}
\begin{points}{对方最可能问的三个问题 → 回应}
\item 爱因斯坦没找到词，人类知不知道？→ 晚一点知道；那一步已经是人走出去的。
\item 说不出的东西，下一代从哪里接？→ 从手上接，从做里接；每一代都有人接到。
\item 你方是不是说语言不重要？→ 不是，语言是人类最好的工具；工具不是边界。
\end{points}
\end{stage}

\begin{kit}{9. 反三申论}{套件 · 组装后 376–393 字 / 90 秒}
\assembly{开头 → 回应正三（任选 1） → 处理正方最强点（任选 1） → 推进论点二 → 收尾}
\begin{fixed}{开头}
谢谢主席。我先回应正三，再处理正方最有力的一点，最后往前推我方第二点。
\end{fixed}
\begin{pick}{回应正三}{1}
\begin{module}{如果正三申论主打性骚扰这个词}
正三刚才讲了“性骚扰”这个词。可这个词出现以前，她们说得清。造词的法利自己写过，那时的人得跟家里解释：他对我动手动脚，我拒绝了还不罢休，我只好辞职。缺的是名字，不是理解；名字，是她们自己造的。
\end{module}
\begin{module}{如果正三申论主打师傅走了说不出的就断了}
正三刚才说，师傅走了，说不出的就断了。可下一个人会再摸到。每一代都有人摸到，这恰恰说明人到得了。说不出，不等于没人到得了；只是每一次，都得人自己去到。所以手艺几千年，没有断过，也没有等过语言。
\end{module}
\begin{module}{如果正三申论主打失语者是先有语言的}
正三刚才说，失语者是先有语言的。可《自然》那篇综述写得清楚：孩子身上，语言和思考早就分开了；没有语言的聋孩子，照样学会数学，照样做因果推理，只是有些慢。你看，慢，从来不是边界；到不了，才是。
\end{module}
\begin{fallback}{正三申论不在以上几条时}
正三刚才那一篇，请评委用判准量一量：它证明的是语言有用，还是语言到不了的地方，人就到不了。前一个，我方从来不反对；后一个，才是今天要辩的。正方要证明的，是后一个。正三讲得再好，也要过这一关。
\end{fallback}
\end{pick}
\begin{pick}{处理正方最强点}{1}
\begin{module}{如果正方全场最有力的是尼加拉瓜手语者}
正方全场最有力的，是尼加拉瓜第一批手语者。我方承认，他们学到那些词以后，答对了。可这说明词能帮人学会，不说明人到不了。那门手语是谁造的？是一群原本没有语言的孩子。人先到了语言外面，才造出了语言，再把后来的人带进去。这个例子，从头到尾都是人走在词前面，词只是跟在后面，把路记了下来。
\end{module}
\begin{module}{如果正方全场最有力的是石器实验说话翻四倍}
正方全场最有力的，是石器实验说话翻四倍。我方承认，说话快。可只比划的那组，也翻了一倍；不用语言，照样教得会，学的人照样学得会。快四倍，是工具好用。好用的工具让人走得快，不划人能走多远。汽车再快，也没人说公路外面就不是人类的地方。何况这个实验里，不出声的那组，也在往前走。
\end{module}
\begin{module}{如果正方全场最有力的是一个人摸到不等于人类走到}
正方全场最有力的，是一个人摸到不等于人类走到。我方承认这句话好听。可人类是谁？就是一个个人。一个人到了，人类的边界就已经在那里；词是后来去接他的，不是先去划他的。把一个人开除出人类，人类就只剩下图书馆了。可图书馆里的每一本书，都是先有一个人想到，才有人写下。
\end{module}
\begin{fallback}{正方最有力的一点不在以上几条时}
我方先把判准再说一遍。正方要证明的，是语言到不了的地方，人就稳定地到不了。我方承认，语言帮人学得快、记得牢。可只要人在语言外面一再到得了，失语的人一次次解出算式，手艺一代代有人摸到，作曲家说不出话还在写曲子，这个“是”字就立不住。这三件事，都没有等语言。
\end{fallback}
\end{pick}
\begin{fixed}{推进论点二}
我再往前推一步。俄语里，浅蓝和深蓝是两个词；英语里，只有一个“蓝”。二〇〇七年，心理学家让二十六个俄国人、二十四个英国人分辨同一组蓝色。俄国人快一些，快了一百二十多毫秒。可准确率呢？两边都在九成六上下。你想想，英国人没有那个词，照样分得清。看得出的，比叫得出的多；做得到的，比说得出的多。
\end{fixed}
\begin{fixed}{收尾}
所以请评委记住：语言让人走得快一点，却从来不划人能走多远。这就是我方第二点。
\end{fixed}
\end{kit}

\begin{stage}{10. 反三被质询防守预案}{正一质询，90 秒双边计时}
\textbf{原则}：俄罗斯蓝的数字报得出出处：Winawer 等，二〇〇七年，美国科学院院刊。对方会说“语言在实时参与感知”，先认，再守“调速度不是划边界”。
\begin{qa}
\qarow{您方说的棋感、手感，说不出来，对吗？}{对。}
\qarow{师傅走了，徒弟没摸到，这一点还在不在？}{在别的摸到的人手里。每一代都有人到得了。}
\qarow{每一代都要从零去摸，算不算往前走？}{每一代都到得了，这就说明边界不在语言那里。}
\qarow{英国人没有两个“蓝”字，也分得清，对吗？}{对，准确率九成六上下。}
\qarow{俄国人更快，一做言语干扰优势就没了，语言是不是在实时参与？}{是，参与的是速度。快一百多毫秒是调速度，分得清分不清才是边界。}
\qarow{分得清颜色，是不是人超出动物的那一截？}{不是。我方用它说明看得出的比叫得出的多；棋感、医生的直觉才是人独有的，也在语言外面。}
\end{qa}
\end{stage}

\begin{stage}{11. 反一质询正三问题链}{90 秒，双边计时}
\textbf{总目标}：为反一结辩取材：拿到“名字是她们自己造的”，和“失语的人还能思考”这个全场问题的正面回答。
\begin{chain}{链一：名字}{人先到，词后到}
\q{一九七五年以前，被骚扰的人知不知道自己受了委屈？}{答“知道”：下一问；答“不知道”：请问她们为什么辞职}
\q{这个名字，是不是她们自己开会想出来的？}{答“是”：下一问}
\q{那是人先到，还是词先到？}{答“人先到”：收束；答“词先到”：词是谁造的}
\close{谢谢。名字是她们自己造的：人先到，词后到。}
\end{chain}
\begin{chain}{链二：全场问题}{失语的人还能思考}
\q{一个中风失语的人，还能不能思考？}{答“能”：下一问；回避：再问一次，记下来}
\q{他的思考，算不算人类的一部分？}{答“算”：收束；答“只算一个人的”：那人类是谁}
\close{谢谢。失语的人还能思考，他的思考也在人类里面。}
\end{chain}
\end{stage}

\begin{stage}{12. 反二对辩要点}{双方各 90 秒}
\textbf{目标}：结辩前最后一次交锋，守住中立判准的“稳定”。让评委看到：我方的例子不是一闪而过的感受，是一再做得到、代代有人到得了的本事。
\begin{points}{进攻点（4–5 个，每个 15–20 秒）}
\item 失语者一次次解出括号算式，这算不算稳定？
\item 手艺几千年没断，这算不算留得住？
\item 正方说一个人摸到不算人类走到。那请问，人类是由谁组成的？
\item 正方的判准要证明“语言到不了，人就到不了”。英国人没有那个词，分蓝色九成六，这是到了还是没到？
\end{points}
\begin{points}{备用点（6–8 个）}
\item 对方说七十万年停滞：停住的是速度，不是人；守住本身就是到得了。
\item 对方说说话翻四倍：快是工具好用。
\item 对方说脚手架：孩子身上早就分开，没语言的聋孩子照样学会数学。
\item 对方说轮扁：不能以喻，是用话教不会。
\item 对方说聋孩子需要语言：我方同意，给孩子语言是好事；好事不等于边界。
\item 对方说维特根斯坦：他自己写，确实有不可说的东西。
\end{points}
\begin{points}{对方最可能问的三个问题 → 回应}
\item 说不出的东西，怎么让下一个人接着往前走？→ 做给他看，手把手，比划也翻一倍。
\item 你们的例子是不是都是一个人？→ 失语研究是一群人，手艺是一代代人。
\item 你们是不是在说语言不重要？→ 语言是最好的工具。工具不是边界。
\end{points}
\end{stage}

\begin{kit}{13. 反一结辩}{套件 · 组装后 379–395 字 / 90 秒}
\assembly{归纳分歧（任选 1） → 处理最强点（任选 1） → 封路（任选 1） → 论点收束 → 价值升华 → 结尾}
\begin{pick}{归纳分歧}{1}
\begin{module}{如果全场主战场落在失语者能不能思考}
谢谢主席。今天吵得最多的，是失语的人还能不能思考。我方的回答是能，研究的回答也是能。语言坏了，思想还在。那么，两条边界就不是同一条。
\end{module}
\begin{module}{如果全场主战场落在说不出的能不能传下去}
谢谢主席。今天吵得最多的，是说不出的东西传不传得下去。手艺几千年没断，比划也教得会。守得住，就是人到得了。那么，两条边界就不是同一条。
\end{module}
\begin{fallback}{主战场不清楚时}
谢谢主席。回到判准：语言到不了的地方，人类还能不能稳定地到达。我方给了两块地方，一块是思考，一块是说不出的本事。两条边界，就在这里分开。
\end{fallback}
\end{pick}
\begin{pick}{处理最强点}{1}
\begin{module}{如果正方全场最有力的是尼加拉瓜手语者}
正方最有力的，是尼加拉瓜的手语者。我方承认，学到那些词，他们答对了。可别忘了，那门手语，是一群原本没有语言的孩子造出来的。人先到了语言外面，才有了语言；词是跟在人后面来的，不是走在人前面的。
\end{module}
\begin{module}{如果正方全场最有力的是石器实验说话翻四倍}
正方最有力的，是说话翻四倍。我方承认，说话快。可只比划的那组，也翻了一倍，不用语言照样教得会。快四倍，是工具好用；好用的工具让人走得快，不是人的边界。不出声的那组，也在往前走。
\end{module}
\begin{fallback}{正方最有力的一点不在以上几条时}
我方承认，语言帮人学得快、记得牢、传得远，这些好处我方一个都不否认。可正方要证明的，是语言到不了的地方，人就到不了。帮了忙，不等于划了线。这一步，是整道题的关键；只要跨不过去，“是”字就立不住。
\end{fallback}
\end{pick}
\begin{pick}{封路}{1}
\begin{module}{如果正方全场主打一个人摸到不等于人类走到}
正方全场都在说，一个人摸到不等于人类走到。待会儿正方的结辩，请评委问一句：人类是谁？如果不是一个个人，那这条边界，到底是谁的边界？一个人先到的地方，算不算人类到过？
\end{module}
\begin{module}{如果正方全场主打没有词就走不到}
正方全场都在说，没有词就走不到。待会儿正方的结辩，请评委问一句：那些词，是谁造出来的？造词的人，在造出词之前，站在哪里？如果站在语言外面，那人就走在语言前面。
\end{module}
\begin{fallback}{正方全场路线不清楚时}
待会儿正方的结辩，请评委拿判准去听：它证明的是语言有用，还是语言到不了的地方，人就一定到不了。只有后一个，才能让“是”字成立；前一个，我方从头到尾都同意。
\end{fallback}
\end{pick}
\begin{fixed}{论点收束}
说到底，我方两点仍然成立：想在说前面，知道的比说的多。语言是人类最好的工具，但再好的工具，也不是边界。
\end{fixed}
\begin{fixed}{价值升华}
最后，请各位想一个人。你外公中风以后，说不出你的名字了。亲戚当着他的面说：他现在什么都不懂了。可你把棋盘摆上，他还是先跳马，还是那样看你一眼。

如果语言的边界就是人的边界，他已经站在人的外面了。我们今天争的，不是语言重不重要。是那些说不出话的人，还算不算一个完整的人。
\end{fixed}
\begin{fixed}{结尾}
语言是人留下的脚印。人，总走在脚印前面。谢谢。
\end{fixed}
\end{kit}
\begin{contingency}
\ifthen{时间不够}{先压“处理最强点”，再压“归纳分歧”；价值升华和结尾不删。}
\ifthen{正方全场都在讲聋孩子需要语言}{封路用兜底，再加一句：给孩子语言是好事，好事不等于边界。}
\end{contingency}
\end{document}
